% 第五章示例：说明讨论部分及论文整体组织的写作方法。
\chapter{讨论与论文组织}

\section{研究结果的解释}

讨论部分应围绕研究目标解释结果，而不是再次重复数据。作者需要分析
结果形成的原因、内在机制和适用条件，并说明结果是否支持最初提出的
研究问题或假设。对于不符合预期的结果，也应如实报告并从样本、测量、
方法假设和外部环境等方面分析可能原因。

\section{与已有研究或实践的比较}

应将主要发现与相关文献、行业标准或已有方案进行比较，说明一致之处、
差异之处及可能原因。比较对象必须具有可比性，不能忽略研究对象、
数据口径、实验条件和评价指标的差异。引用文献是为论证提供依据，
而不是用文献数量代替分析\cite{kanamori1998shaking,zhu1996fluid}。

\section{研究贡献与局限}

研究贡献应从实际完成的工作和证据出发，可以体现在新的研究对象、
数据资料、分析视角、方法改进、系统实现或应用验证等方面。贡献表述
需要与摘要和总结一致，并明确哪些成果由作者本人完成。

任何研究都受到样本、时间、数据质量、实验条件和方法假设等限制。
主动说明局限不会削弱论文价值，反而有助于读者正确理解结论边界。
局限之后可提出具体的改进方向，但不应把本应完成而未完成的核心工作
简单留作展望。

\section{论文整体组织与证据链}

\subsection{标题和章节关系}

各级标题应简明扼要，准确概括所辖内容。相邻章节之间要有清楚的逻辑
过渡，避免同一内容在多章重复出现。绪论提出的问题、研究方法、结果、
讨论和总结应形成完整证据链，最终结论必须能够在前文找到依据。

\subsection{图表、公式和正文引用}

图表与公式应先在正文中提及，再给出具体内容。例如，本示例中的
图~\ref{fig:distribution}、表~\ref{tab:gear} 和式~\eqref{eq:fourier}
均在相关段落中得到解释。修改章节或移动内容后，编号和目录应统一
更新，避免出现正文所述编号与实际对象不一致的情况。

\subsection{目录与全文复核}

目录是论文结构的集中呈现，通常列到三级标题，并应与正文标题完全
一致。定稿前应从头到尾检查题目、摘要、目录、图表清单、正文、
参考文献、致谢和附录，重点核对术语、符号、单位、编号、引用和结论
是否前后一致。

\section{学术诚信与资料留存}

使用他人的观点、文字、数据、图片、程序或方法时，应按照规范明确
标注来源。实验记录、调查问卷、数据处理过程、程序版本和重要中间
结果应妥善保存，以便指导教师审阅和后续复核。不得伪造、篡改数据，
也不得通过改写措辞掩盖未经标注的复制使用。
